\begin{titlepage}
  \centering
  \vspace*{2cm}
  {\LARGE\bfseries Дискретная информационная\\[0.4em] термодинамика Вселенной\par}
  \vspace{1.2cm}
  {\large Монография о решёточной жидкости на планковском масштабе:\\
  от локального закона эволюции к наблюдаемой физике\par}
  \vfill
  {\large \today\par}
\end{titlepage}

\chapter*{Предисловие}
\addcontentsline{toc}{chapter}{Предисловие}

Настоящая монография строится от первых принципов, по порядку введения объектов.
Сначала --- короткий вопрос: что уже проверено в макрофизике и где этот язык упирается в предел.
Затем вводится носитель как место для состояния: узлы, шаг, такт, окрестность, поле на узле.
Только на этом языке формулируются абсолютные условия.
Геометрия графа, закон эволюции и вывод макроописания идут после них,
а не до того, как назван объект, о котором они говорят.

Изложение --- как в теоретическом учебнике: определения, теоремы и доказательства
связаны перекрёстными ссылками.
Краткие пояснения по обратимым дискретным схемам и представлению Маделунга ---
в приложении~\ref{app:background}.

Читатель, знакомый с квантовой механикой, термодинамикой и специальной теорией относительности,
может идти по тексту последовательно: объект, условия, теоремы, восстановление известной физики.

\tableofcontents
